\section{The Compactness Distribution of ReCom Ensembles}
\label{sec:distribution}

\subsection{Right-Skewed Shape}

The compactness distribution of \recom{} ensembles is typically right-skewed:
most plans cluster at moderate compactness values, with a long right tail of
unusually compact plans and a shorter left tail of very non-compact plans.

This shape arises from the recombination move structure.  A \recom{} step
merges two adjacent districts and redraws the boundary using a random spanning
tree cut.  Random spanning tree cuts tend to produce boundaries with many
turns and jags, which reduces Polsby-Popper compactness.  Plans with
high compactness require specific boundary choices that the random process
rarely achieves.  Hence, the distribution has a mode well below the maximum
possible compactness.

Formally, if $X \sim F_{\pp}$ is the mean Polsby-Popper score across the
ensemble, we observe empirically:

\begin{equation}
  \mathrm{Skewness}(F_{\pp}) \approx +0.3 \text{ to } +0.7
  \quad\text{(positive, right-skewed)}
  \label{eq:skew}
\end{equation}

for the six states studied.  The skewness is higher for geographically
irregular states (Pennsylvania, California) and lower for more regular states
(Wisconsin).

\subsection{Typical Values}

Typical mean Polsby-Popper values in \recom{} ensembles are:

\begin{center}
\begin{tabular}{lcccc}
\toprule
State & $k$ & Ensemble mean & Ensemble 75th pct & Ensemble max \\
\midrule
NC & 14 & 0.318 & 0.355 & $\approx 0.52$ \\
WI &  8 & 0.301 & 0.342 & $\approx 0.49$ \\
GA & 14 & 0.312 & 0.348 & $\approx 0.51$ \\
PA & 17 & 0.306 & 0.341 & $\approx 0.48$ \\
\bottomrule
\end{tabular}
\end{center}

The \ar{} plan falls between the 75th percentile and the ensemble maximum
for all states: it is more compact than the average \recom{} plan but
well below the maximum.

\subsection{Why the Maximum Is Not Achieved}

The \recom{} ensemble maximum (plans with $\bar\pp \approx 0.49$--$0.52$)
represents extremely compact plans that exist mathematically but rarely appear
in the chain.  These plans typically have districts that align with
major geographic features (rivers, mountain ranges) and avoid coastal
meanders and complex inland topographies.

The \ar{} plan does not target Polsby-Popper directly.  It targets
METIS edge cut, which is related to but not identical to Polsby-Popper.
A district with minimum edge cut will have few cross-district census-tract
adjacencies, which correlates with a smooth boundary, which correlates with
high Polsby-Popper.  But the correlation is not perfect: some minimum-edge-cut
plans have irregular boundaries that happen to avoid many cross-boundary
adjacencies due to the geographic structure of the tract network.

This distinction explains why the \ar{} plan is at the $61$st--$72$nd
percentile rather than the 99th percentile: it is compact, but not
maximally compact in the Polsby-Popper sense.
